% Modernised reading — fonds Alexandre Grothendieck, shelfmark 11, pages 1-57.
%
% This is an INTERPRETATION, not a transcription. It re-reads the pages in
% today's notation and vocabulary, states the hypotheses the manuscript leaves
% implicit, and drops the critical apparatus so that the mathematics reads
% straight through. Everything the brackets and underlines carried moves into a
% footnote. It is held to being mathematically correct as it stands; where the
% manuscript is loose or mistaken, the true statement is what appears here and
% the footnote says what the page has. In any dispute of substance the
% transcription governs: cite batch-01/02/03.fr.tex.
%
% Pass: Opus 5 (claude-opus-5), 2026-09-14 — first pass, unchecked against the
% pages by a human.
%
% Scope: one document for the shelfmark, its three batches read as one
% argument. Only the even leaves 2-30, then 33-39 odd, 41-47 odd and 50-56 even
% carry his hand; the rest is the back of an unrelated English typescript, a
% blank page (1, 32) or show-through (49), and is not mentioned.
%
% Spine. (1) pp. 2-18: from a « bivariant » functor A on a category V with
% products (pull-backs, push-forwards, projection and exchange formulas) he
% builds the category of correspondences V^A, Hom(X,Y) = A(X×Y), with the
% graph functor φ, the action ψ, transposition, and a remark that only the
% monoid structure was used. (2) pp. 20-24: sub-functors A' ⊂ A stable under
% the operations; their criterion; the smallest one A_sp; generation by maps
% « of type A' ». (3) pp. 26-35: Poincaré algebras and then algebras with a
% form, where f_* is determined by f^*; the sign of (f×g)_*; the degree shift
% d(B)-d(A). (4) pp. 37-39: the tensor product is not a product; simplicial and
% « special » bivariant algebras; a Proposition characterising stable systems
% A_n ⊂ H*(X^n). (5) pp. 41-47: a SECOND MOUTURE of the same programme in the
% simplicial language (gauged algebras, examples, correspondences between
% powers X^I, units Δ_I), breaking off before associativity. (6) pp. 50-56: the
% tautological sub-algebra of H*(X^n), generated by the point class y and the
% diagonal classes. The two formalisms of correspondences (pp. 2-16 and 43-47)
% are kept distinct: the second changes the order of the factors and forgets
% the exchange axiom, and he notes both.
%
% NOTATION fixed for the folder, with the letters he reuses renamed:
%   V — base category (p. 2); but « V° » on p. 26 is the category of Poincaré
%     algebras, written P here, and « V » on p. 41 is a category of varieties,
%     written Var here.
%   C — category of bivariant algebras (pp. 2, 37); the element C ∈ A(I⊔I') of
%     p. 43 is written c.
%   φ — the graph functor V → V^A (p. 8) only. The functor f^* ↦ f_* of p. 26
%     is written f^* ↦ f_* with no letter; the map A → Ǎ of p. 33 is θ_A.
%   ψ(u) — written ũ (his notation from p. 16 on).
%   κ_A — the form (his k_A, k_X).
%   projections — p_{12}, p_{23}, p_{13} for X×Y×Z → X×Y etc., everywhere;
%     his p_{21}, p_{32}, p_{31} are these.
%   composition — v∘u = p_{13*}(p_{23}^*(v)·p_{12}^*(u)) (p. 2's order);
%     p. 45 takes the other order, footnoted.
%   δ_I — class of the small diagonal of X^I (δ_{{i}} = 1); p. 52's δ_r is
%     δ_I for |I| = r+1.
%   γ_f = Γ_{f*}(1) — class of the graph (p. 20).
%
% Standing hypothesis supplied: all A(X) graded-commutative and all push-
% forwards f_* of even degree (p. 30 adopts it from there on; p. 35 names the
% even subcategory). Under it every identity of pp. 2-24 holds without sign; the
% pages mark « att. signes ! » throughout and never fix the signs otherwise.
%
% Substantive departures, each footnoted where it occurs:
%   - p. 2: the exchange axiom is required for a class of transversal cartesian
%     squares (those p. 8 lists), not for all cartesian squares — for the
%     square of two diagonals it fails by the Euler class, which p. 50's
%     δ² = χ y_1 y_2 records;
%   - p. 14: the computation ends « = x » where it gives f_*(x); « f : Y → X »
%     where X → Y is needed; s written V^{A°} → V^{A°} where (V^A)° → V^A;
%   - p. 22: A_sp(X) is the GROUP generated by the f_*(1_Y), not their set (the
%     set is not closed under products: E·E = −[pt] on a blow-up);
%   - p. 24: B(X) must be the set of finite SUMS of values of maps of type A';
%   - p. 28: the boxed sign formula is legible only in part; recomputed under
%     explicit conventions, agreeing with the page that no sign appears when
%     deg f is even;
%   - p. 39 (i) a): u(x) ∈ A_m, where the page writes H*(X^m);
%   - p. 41: degrees in [0, d(A)], where the page reads [1, d(A)];
%   - p. 41: full faithfulness of A(·) ↦ (A(pt), Δ) is clear only for special
%     objects;
%   - p. 50: δ ∈ H^n(X×X), where the page writes H^n(X); the « bases » are
%     spanning families, free only generically;
%   - pp. 54-56: the product of two classes ε_{R,J} is restated as a complete
%     rule; the coefficient in the case s + s' = k, which the page leaves with a
%     struck χ and a « ? », is χ by the excess intersection formula — checked
%     against p. 50's two instances; the pointed-block case is ours;
%   - p. 56: the push-forward of δ_I along a projection is δ_{I∖{i}}, not δ_I;
%     « δ_k » for the small diagonal of X^k is δ_{[1,k]}.
%
% Announced and never established, recorded in the spine section: the precise
% exchange axiom (p. 2; « Question » p. 30; « axiome cartésien oublié ! »
% p. 45); multiplicativity of s (p. 14, « vérifier ») and ψφ* = ω*A° (p. 16,
% « à vérifier »); the theory of graded monoids with involutions (p. 18); the
% more general statement with a set H of morphisms (p. 20, margin); what A_sp
% is generated by (p. 22, the line stops); the equivalence of special
% simplicial algebras with Poincaré algebras (pp. 37, 41, no proof, « à
% vérifier ») and the axioms on Δ (p. 41); a proof of the Proposition of p. 39;
% associativity and units in the second mouture (p. 47 stops on a comma); the
% relations of H*(X^k) (p. 52, heading only); the d) of p. 56 (struck).
%
% Modern names are footnoted with whether the pages could have known them.
% Nothing on the leaves names Manin, Kleiman, Demazure, motives, Green functors
% or Frobenius algebras, and this reading does not imply otherwise; the dating
% [à partir de 1961] is the archivists'.
\documentclass[11pt,a4paper]{article}
% Le préambule commun à toutes les transcriptions du fonds Grothendieck.
%
% Il tient en une idée : **l'incertitude se compose, elle ne se résout pas.**
% Une transcription qui lisse un mot douteux a détruit la seule chose pour
% laquelle on la faisait — un lecteur doit pouvoir savoir, à chaque endroit,
% ce qui était sur le feuillet et ce qui a été ajouté. Les six macros qui
% suivent sont donc l'appareil critique, et non de la décoration.
%
% Les mêmes commandes sont reconnues par `scripts/render.mjs`, qui produit la
% vue de lecture du site. Une macro ajoutée ici doit l'être aussi là-bas,
% faute de quoi le rendu échouera bruyamment — ce qui est voulu : un rendu
% silencieusement faux est pire qu'un rendu absent.

\ProvidesPackage{grothendieck}[2026/08/08 Transcriptions du fonds Grothendieck]

\usepackage{amsmath,amssymb,amsthm}
\usepackage{mathrsfs}
\usepackage{stmaryrd}
% Les diagrammes commutatifs sont partout dans ce fonds : c'est la langue
% même de la géométrie algébrique de Grothendieck, et une transcription qui
% les rendrait en prose ne transcrirait rien.
\usepackage{tikz-cd}
% Français par défaut — c'est la langue des feuillets — mais l'anglais est
% chargé aussi : la traduction et le résumé sont des documents anglais, et la
% typographie française (espaces avant les deux-points) y serait une faute.
% Ces documents basculent par \selectlanguage{english} en tête de corps.
\usepackage[english,french]{babel}
\usepackage{fontspec}
\usepackage{geometry}
\geometry{a4paper,margin=25mm,marginparwidth=28mm}
\usepackage{marginnote}
\usepackage{xcolor}
% ulem, not soul. `soul`'s \st and \ul are built on a character-by-character
% scan that XeLaTeX's font machinery breaks: the document fails with an
% undefined control sequence rather than degrading. ulem does the same two jobs
% and is Unicode-engine safe. `normalem` stops it hijacking \emph, which the
% transcriptions use for Grothendieck's own emphasis.
\usepackage[normalem]{ulem}
\usepackage{hyperref}
\hypersetup{colorlinks=true,linkcolor=black,urlcolor=black,pdfborder={0 0 0}}

% --- Métadonnées du lot -----------------------------------------------------
% Reprises telles quelles par le script de rendu, qui les affiche en tête de
% la vue de lecture. Elles fixent ce que le fichier prétend couvrir : sans
% elles, une transcription flotte sans point d'attache dans un fonds de seize
% mille feuillets.

\newcommand{\folder}[1]{\gdef\grothFolder{#1}}
\newcommand{\batch}[1]{\gdef\grothBatch{#1}}
\newcommand{\leaves}[2]{\gdef\grothFirst{#1}\gdef\grothLast{#2}}
\newcommand{\foldertitle}[1]{\gdef\grothFoldertitle{#1}}
\gdef\grothFolder{?}\gdef\grothBatch{?}\gdef\grothFirst{?}\gdef\grothLast{?}\gdef\grothFoldertitle{}

% --- Le résumé --------------------------------------------------------------
%
% Il ouvre la lecture modernisée, sous le titre « Résumé », et il est composé
% en petit. Ce n'est pas un ornement : c'est le seul passage du document qui
% ne soit pas des mathématiques, et un lecteur doit pouvoir le reconnaître
% avant de l'avoir lu — puis le sauter s'il connaît déjà le sujet, ou s'y
% arrêter s'il ne le connaît pas. Le filet qui le suit marque la reprise du
% corps.

\newenvironment{resume}
  {\small\setlength{\parskip}{0.45em}}
  {\par\medskip\hrule\medskip}

% --- Mots-clés ---------------------------------------------------------------
%
% En anglais, en clôture du résumé de la lecture modernisée : le vocabulaire
% moderne sous lequel chercher ce que les feuillets construisent. Le manifeste
% du site (`npm run manifest`) les extrait pour étiqueter la cote entière —
% cette ligne est la source unique des tags, il n'y a pas de fichier à côté.
\newcommand{\keywords}[1]{\par\smallskip\noindent
  {\footnotesize\textsc{Keywords} --- \textit{#1}}\par}

% --- L'appareil critique ----------------------------------------------------

% Le numéro de feuillet, en marge. C'est l'ancre qui permet de régler un
% désaccord sur une formule en regardant *un* feuillet plutôt qu'une chemise
% de six cents. Le site s'en sert aussi pour tourner les pages du fac-similé
% quand on fait défiler la transcription.
\newcommand{\leaf}[1]{%
  \marginnote{\footnotesize\color{gray!70}\textsf{#1}}%
  \label{leaf:#1}\ignorespaces}

% Illisible : on ne devine pas. Un mot inventé qui se lit comme les autres est
% le pire résultat possible.
\newcommand{\ill}{\textcolor{red!60!black}{[\,\ldots\,]}}

% Lecture douteuse : le mot est donné, et signalé comme incertain.
\newcommand{\uncertain}[1]{\uline{#1}}

% Ajout de l'éditeur — une lettre manquante, un mot sous-entendu.
\newcommand{\add}[1]{\textcolor{blue!50!black}{[#1]}}

% Biffé par Grothendieck lui-même. Ce qu'il a rayé fait partie du document :
% une rature dit souvent où la pensée a changé de direction.
\newcommand{\struck}[1]{\sout{#1}}

% Note du transcripteur, sur l'état matériel du feuillet.
\newcommand{\note}[1]{\par\smallskip{\small\color{gray!60!black}\textit{[#1]}}\par\smallskip}

% Note marginale de Grothendieck, distincte du corps du texte.
\newcommand{\marginal}[1]{\par\smallskip\noindent\hspace{1em}%
  {\small\color{gray!60!black}#1}\par\smallskip}

% --- Titre ------------------------------------------------------------------

\AtBeginDocument{%
  \begin{center}
    {\large\bfseries Fonds Alexandre Grothendieck --- cote n\textsuperscript{o}~\grothFolder}\\[2pt]
    {\itshape\grothFoldertitle}\\[4pt]
    {\small feuillets \grothFirst--\grothLast\ (lot \grothBatch)}\\[2pt]
    {\footnotesize\color{gray!60!black}Transcription établie d'après le fac-similé numérisé
     par l'Université de Montpellier.\\
     Les lectures incertaines sont soulignées, les illisibles notées [\,\ldots\,],
     les ajouts entre crochets.}
  \end{center}
  \vspace{1em}}

\endinput


\folder{11}
\pages{1}{57}
\dating{[à partir de 1961]}
\watermark{Édition de démonstration\\interprétation personnelle de l'œuvre}
\foldertitle{Formalisme des correspondances : notes manuscrites (s.d.) — lecture modernisée du dossier entier}

\begin{document}

\section*{Résumé}

\begin{resume}
Ces feuillets cherchent les règles les plus économes qui permettent de
\emph{composer des correspondances} comme on compose des applications : de
fabriquer, à partir de la cohomologie des espaces, une catégorie dont les
flèches sont des classes vivant sur des produits.

L'image la plus simple est celle des matrices. Une matrice à lignes indexées
par un ensemble fini $X$ et colonnes indexées par $Y$ est une fonction sur le
produit $X \times Y$. Pour composer deux matrices, on les prolonge toutes deux
à $X \times Y \times Z$, on les multiplie point par point, puis on somme sur
l'indice du milieu. La matrice identité est la fonction qui vaut $1$ sur la
diagonale et $0$ ailleurs ; une application $f \colon X \to Y$ donne la matrice
de son graphe ; une matrice agit sur les fonctions sur $X$ par la même
recette. Rien de tout cela n'utilise autre chose que trois opérations : tirer
une fonction en arrière le long d'une application, multiplier, et sommer sur
les fibres.

L'idée du dossier est que la cohomologie d'un espace offre exactement ces
trois opérations : l'image inverse $f^{*}$, le cup-produit, et l'image directe
$f_{*}$, qui « intègre sur les fibres ». Deux formules les relient, et ce sont
les seules qu'on utilise : la formule de projection, et une formule
d'échange, qui dit qu'on peut permuter image inverse et image directe dans un
carré cartésien. Avec elles, les pages 2 à 16 construisent la catégorie des
correspondances : ses objets sont les espaces, les flèches de $X$ vers $Y$ sont
les classes de cohomologie de $X \times Y$, et la composition est la recette
des matrices. Chaque fois qu'il a fini une vérification, il fait le compte de
ce dont il s'est servi --- tel cas de la formule d'échange et pas tel autre, la
multiplication et jamais l'addition --- et c'est ce compte qui dessine les
axiomes. Les pages 20 à 24 cherchent ensuite les sous-familles de classes qui
sont stables par toutes ces opérations ; la plus petite, dans les cas
géométriques, est formée des classes de sous-variétés.

La seconde moitié change de point de vue. Pour une algèbre de cohomologie qui
satisfait la dualité de Poincaré, l'image directe est entièrement déterminée
par l'image inverse ; les pages 26 à 39 font de ces algèbres les objets d'une
catégorie, découvrent au passage que les signes s'y dérèglent dès que les
degrés sont impairs, et repèrent qu'on ne peut pas se borner à ces algèbres-là.
Les pages 41 à 47 reprennent tout le programme dans un langage où les espaces
sont remplacés par les puissances $X^{I}$ d'un seul espace.

La fin du dossier (pages 50 à 56) calcule. Dans la cohomologie des puissances
$X \times \cdots \times X$, on ne garde que ce qu'on peut fabriquer à partir de
la classe d'un point et des diagonales, et on cherche la table de
multiplication de ces classes. Quand deux diagonales se coupent « comme il
faut », leur produit est une diagonale ; quand elles se coupent mal, il
apparaît un nombre, la caractéristique d'Euler de $X$. C'est l'analogue exact de
la trace de la matrice identité, qui compte les points de $X$ : pour un espace,
la diagonale coupée par elle-même compte, avec leurs signes, les points fixes
de l'identité. La page 30 avait demandé dans quels cas la formule d'échange
reste vraie ; ce calcul en donne la réponse géométrique, sans que le dossier la
formule comme telle.

C'est le socle de ce qu'on appelle aujourd'hui les motifs purs : la catégorie
des correspondances d'une cohomologie, son sous-anneau de classes algébriques,
ses relations d'équivalence (numérique, homologique). On la retrouvera aussi
sous les noms de foncteur de Green, d'algèbre de Frobenius, de formule
d'intersection en excès et d'anneau tautologique.

\keywords{category of correspondences, pure motives, Weil cohomology, projection formula, base change, Green functor, algebraic cycle class, Chow ring, adequate equivalence relation, Poincaré duality algebra, Frobenius algebra, Koszul sign rule, symmetric simplicial object, diagonal class, excess intersection formula, Euler characteristic, tautological ring}
\end{resume}

\pagerange{2}{2}
\section*{Le fil du dossier, et les conventions}

\emph{Les stations.} Le dossier suit un seul programme, mené deux fois.

\begin{enumerate}
  \item \emph{Pages 2 à 18.} Une théorie « bivariante » $A$ sur une catégorie
  $\mathcal{V}$ à produits ; la catégorie des correspondances
  $\mathcal{V}^{A}$ ; le foncteur des graphes, l'action sur les $A(X)$, la
  transposition ; une remarque sur ce que les calculs ont réellement utilisé.
  \item \emph{Pages 20 à 24.} Les sous-foncteurs stables $A' \subset A$ :
  critère, plus petit sous-foncteur stable, engendrement.
  \item \emph{Pages 26 à 35.} Les algèbres de Poincaré, puis les algèbres
  munies d'une forme, où l'image directe est déterminée par l'image inverse ;
  le signe de $(f \times g)_{*}$ ; le décalage des degrés.
  \item \emph{Pages 37 et 39.} Le produit tensoriel n'est pas un produit ; les
  algèbres bivariantes simpliciales et spéciales ; une Proposition sur les
  systèmes stables $A_{n} \subset H^{*}(X^{n})$.
  \item \emph{Pages 41 à 47.} Seconde mouture : algèbres jaugées simpliciales,
  exemples, correspondances entre puissances $X^{I}$, unités. Elle s'arrête
  avant l'associativité.
  \item \emph{Pages 50 à 56.} La sous-algèbre tautologique de
  $H^{*}(X^{n})$.
\end{enumerate}

Les deux formalismes des correspondances --- pages 2 à 16 et pages 43 à 47 ---
sont deux \emph{moutures} du même texte et sont gardés distincts : la seconde
change l'ordre des facteurs de la composition, travaille sur les puissances
d'un seul espace, et oublie d'abord l'axiome d'échange, ce que la page 45 note
elle-même.

\emph{Hypothèse de signe.} Sauf mention contraire, les algèbres $A(X)$ sont
commutatives au sens gradué et toutes les images directes $f_{*}$ sont de
degré pair. Sous cette hypothèse, toutes les identités des pages 2 à 24
valent sans signe.\footnote{Les pages 4, 6, 12 et 16 marquent « attention
signes ! » à chaque usage de la formule de projection, et la page 18 envisage
de refaire la théorie pour les contrôler ; le dossier ne fixe jamais ces
signes. C'est la page 30 qui adopte l'hypothèse « de degré pair », et la page
35 qui nomme la sous-catégorie correspondante. Elle couvre les cas qui motivent
le dossier : cohomologie de Betti ou $\ell$-adique des variétés algébriques
lisses, où $f_{*}$ est de degré $2(\dim Y - \dim X)$, anneaux de Chow.}

\emph{Notations, fixées pour tout le dossier.} Plusieurs lettres servent deux
ou trois fois sur les pages ; on les sépare ainsi.

\begin{itemize}
  \item $\mathcal{V}$ est la catégorie de base des pages 2 à 24. La catégorie
  des algèbres de Poincaré, que la page 26 note $\mathcal{V}^{\circ}$, est
  notée ici $\mathcal{P}$ ; la catégorie de variétés de la page 41 est notée
  $\mathrm{Var}$.
  \item $\mathcal{C}$ est la catégorie des algèbres bivariantes ; la classe que
  la page 43 appelle $C$ est notée $c$.
  \item $\varphi$ désigne seulement le foncteur des graphes. L'application
  $A \to \check{A}$ de la page 33 est notée $\theta_{A}$.
  \item L'action d'une correspondance $u$ est notée $\tilde{u}$, et sa
  transposée ${}^{t}u$.
  \item La forme d'une algèbre est $\kappa_{A}$ (il écrit $k_{A}$), pour la
  séparer de l'anneau de base $k$.
  \item Sur $X \times Y \times Z$, $p_{12}$, $p_{23}$, $p_{13}$ sont les
  projections sur $X \times Y$, $Y \times Z$, $X \times Z$.\footnote{Il écrit
  $p_{21}$, $p_{32}$, $p_{31}$ dans les formules et $p_{12}$, $p_{23}$,
  $p_{13}$ sur les diagrammes de la même page 2 ; ce sont les mêmes flèches.}
  \item $\gamma_{f} = \Gamma_{f*}(1)$ est la classe du graphe de $f$ ;
  $\Delta_{X} = \gamma_{\mathrm{id}_{X}}$ celle de la diagonale ; $\delta_{I}$
  la classe de la petite diagonale de $X^{I}$, avec $\delta_{\{i\}} = 1$.
\end{itemize}

\emph{Ce que le dossier annonce sans l'établir.} L'énoncé exact de l'axiome
d'échange (page 2 ; « Question » de la page 30 ; « axiome cartésien oublié ! »
à la page 45) ; la multiplicativité de la transposition (page 14) et
l'identification du côté contravariant (page 16), toutes deux marquées « à
vérifier » ; la théorie des monoïdes gradués à involutions (page 18) ; un
énoncé plus général où l'on ne demande la stabilité que pour une famille de
morphismes (marge de la page 20) ; ce par quoi $A_{\mathrm{sp}}$ est engendré
(page 22, la ligne s'arrête) ; l'équivalence entre algèbres simpliciales
spéciales et algèbres de Poincaré (pages 37 et 41, « à vérifier ») et les
axiomes à imposer à $\Delta$ (page 41) ; une démonstration de la Proposition
de la page 39 ; l'associativité et les unités de la seconde mouture (la page
47 s'arrête sur une virgule) ; les relations de $H^{*}(X^{k})$ (page 52, le
titre seul) ; le d) de la page 56, biffé.

\pagerange{2}{2}
\section*{La catégorie des correspondances d'une théorie bivariante (pages 2 à 18)}

\subsection*{Les données}

Soit $\mathcal{V}$ une catégorie possédant les produits finis, donc un objet
final $e$. On se donne un foncteur $A$ de $\mathcal{V}$ dans la catégorie
$\mathcal{C}$ des algèbres graduées bivariantes : pour tout objet $X$, une
$k$-algèbre graduée unitaire $A(X)$ ; pour tout $f \colon X \to Y$, un
homomorphisme d'algèbres unitaires $f^{*} \colon A(Y) \to A(X)$ et une
application $k$-linéaire homogène $f_{*} \colon A(X) \to A(Y)$, avec
$(gf)^{*} = f^{*} g^{*}$ et $(gf)_{*} = g_{*} f_{*}$. Un morphisme de
$\mathcal{C}$ est un tel couple $(f_{*}, f^{*})$, et
$\omega \colon \mathcal{C} \to k\text{-Mod}$ le foncteur qui retient
$f_{*}$.\footnote{Dans $k\text{-Mod}$, les flèches sont les applications
$k$-linéaires entre modules gradués, non nécessairement homogènes : une
correspondance n'a pas de raison d'être homogène. La page 2 écrit « $k$-mod
gr. ».} On demande :

\begin{enumerate}
  \item[(P)] la \emph{formule de projection}
  $f_{*}\bigl(f^{*}(y)\, x\bigr) = y\, f_{*}(x)$ ;
  \item[(E)] la \emph{formule d'échange} $g^{*} f_{*} = f'_{*}\, g'^{*}$ pour
  tout carré cartésien
\end{enumerate}

\begin{tikzcd}
X' \arrow[r, "g'"] \arrow[d, "f'"'] & X \arrow[d, "f"] \\
Y' \arrow[r, "g"'] & Y
\end{tikzcd}

d'une classe de carrés dits \emph{transverses}, qui contient au moins les
carrés formés de projections de produits, et ceux où une diagonale
$\delta_{X} \colon X \to X \times X$ ou un graphe
$\Gamma_{f} = (\mathrm{id}_{X}, f) \colon X \to X \times Y$ est changé de base
par une projection ou par un autre graphe.\footnote{La page 2 demande
seulement « une condition d'échange », sans l'énoncer. C'est la page 8 qui
recense les deux sortes de carrés dont les vérifications se sont servies ; la
page 10 ajoute le carré de deux graphes. Ces carrés sont cartésiens dans
toute catégorie à produits, et transverses dans les exemples géométriques.
Exiger (E) pour \emph{tous} les carrés cartésiens serait faux en cohomologie :
le carré formé de la diagonale $\delta \colon X \to X \times X$ et
d'elle-même est cartésien, mais $\delta^{*}\delta_{*}(1)$ est la classe d'Euler
de $X$ et non $1$ --- c'est ce qu'enregistre la relation
$\delta^{2} = \chi\, y_{1} y_{2}$ de la page 50. Le dossier ne tranche pas :
la page 30 pose la question, et la page 45 note qu'il a « oublié » l'axiome. La
réponse moderne est la formule d'intersection en excès (Fulton,
\emph{Intersection theory}, 1984, § 6.3), qui corrige (E) par une classe
d'Euler quand le carré n'est pas transverse.}

Le modèle est celui des ensembles finis : $A(X) = k^{X}$, $f^{*}$ la
composition avec $f$ et $f_{*}$ la somme sur les fibres. Là, (P) et (E) valent
pour tous les carrés cartésiens, et tout ce qui suit redonne le calcul
matriciel.\footnote{L'exemple est de nous ; le dossier n'en donne aucun avant
la page 41. La structure (P) + (E) sur une catégorie à produits fibrés est ce
qu'on appelle depuis Green (1971) et Dress (1973) un \emph{foncteur de Green} :
un foncteur de Mackey muni d'une structure d'anneau compatible, la formule de
projection y portant le nom de réciprocité de Frobenius. Le mot
« bivariant » des pages n'est pas celui de Fulton et MacPherson (1981), dont
les théories associent un groupe à chaque morphisme.}

\subsection*{La composition}

On pose $\operatorname{Ob} \mathcal{V}^{A} = \operatorname{Ob} \mathcal{V}$ et
\[ \operatorname{Hom}_{\mathcal{V}^{A}}(X, Y) = A(X \times Y) . \]
Pour $u \in A(X \times Y)$ et $v \in A(Y \times Z)$, on définit
\[ v \circ u = p_{13*}\bigl( p_{23}^{*}(v)\, p_{12}^{*}(u) \bigr) \in A(X \times Z) . \]

\begin{tikzcd}
 & X \times Y \times Z \arrow[dl, "p_{12}"'] \arrow[dr, "p_{23}"] \arrow[dd, "p_{13}"] & \\
X \times Y & & Y \times Z \\
 & X \times Z &
\end{tikzcd}

C'est une application $k$-bilinéaire, de sorte que $\mathcal{V}^{A}$ sera une
catégorie $k$-linéaire.

\pagerange{4}{4}
\subsection*{Associativité}

Pour $X \xrightarrow{u} Y \xrightarrow{v} Z \xrightarrow{w} T$, les deux
composés $w \circ (v \circ u)$ et $(w \circ v) \circ u$ valent
\[ p_{XT*}\bigl( p_{ZT}^{*}(w)\, p_{YZ}^{*}(v)\, p_{XY}^{*}(u) \bigr) , \]
où les $p$ sont les projections de $X \times Y \times Z \times T$. Pour le
premier, on écrit $v \circ u$ comme image directe par
$X \times Y \times Z \to X \times Z$, on la remonte sur
$X \times Z \times T$ par la formule d'échange appliquée au carré de
projections
\[ X \times Y \times Z \times T \to X \times Z \times T , \qquad X \times Y \times Z \to X \times Z , \]
puis on fait entrer $p_{ZT}^{*}(w)$ sous l'image directe par la formule de
projection et l'on compose les images directes. Le second se traite de même,
en passant par $X \times Y \times T$ et $Y \times Z \times T$. Seuls servent
des carrés de projections.

\pagerange{6}{8}
\subsection*{Unités}

Soit $\Delta_{X} = \delta_{X*}(1) \in A(X \times X)$. On a
$u \circ \Delta_{X} = u$ pour $u \in A(X \times Y)$. En effet, le carré

\begin{tikzcd}
X \times Y \arrow[r, "q = \delta_{X} \times \mathrm{id}_{Y}"] \arrow[d, "p_{1}"'] & X \times X \times Y \arrow[d, "p_{12}"] \\
X \arrow[r, "\delta_{X}"'] & X \times X
\end{tikzcd}

est cartésien, d'où $p_{12}^{*}(\Delta_{X}) = q_{*}(1)$ par échange ; puis
$p_{23}^{*}(u)\, q_{*}(1) = q_{*}\bigl(q^{*} p_{23}^{*}(u)\bigr) = q_{*}(u)$
par projection, puisque $p_{23}\, q = \mathrm{id}_{X \times Y}$ ; enfin
$u \circ \Delta_{X} = p_{13*}\, q_{*}(u) = u$, puisque
$p_{13}\, q = \mathrm{id}_{X \times Y}$. On montre de même
$\Delta_{Y} \circ u = u$.\footnote{La page 6 ne fait que le premier calcul et
dit que l'autre est analogue ; une note en marge, en partie douteuse, renvoie
aux signes $(-1)^{ux}$ de ce second calcul.} Ainsi $\mathcal{V}^{A}$ est une
catégorie.

Il fait alors le compte de ce qui a servi. La composition et l'associativité
n'emploient que les projections et les carrés de projections ; les unités
exigent en plus les diagonales, et le carré où une diagonale est changée de
base par une projection. Cette remarque est ce qui reviendra, page 43, sous la
forme « il suffit de connaître les opérations face ».

\pagerange{8}{12}
\subsection*{Le foncteur des graphes et l'action}

Pour $f \colon X \to Y$ dans $\mathcal{V}$, on pose
\[ \varphi(f) = \gamma_{f} = \Gamma_{f*}(1) \in A(X \times Y) . \]
Alors $\varphi(\mathrm{id}_{X}) = \Delta_{X}$ par définition, et
$\varphi(g) \circ \varphi(f) = \varphi(gf)$ pour
$X \xrightarrow{f} Y \xrightarrow{g} Z$. On calcule
$p_{23}^{*}(\gamma_{g})\, p_{12}^{*}(\gamma_{f})$ sur $X \times Y \times Z$ :
chaque facteur est, par échange, l'image directe de $1$ par un plongement
$\mathrm{id}_{X} \times \Gamma_{g}$ ou $\Gamma_{f} \times \mathrm{id}_{Z}$ ;
le carré de ces deux plongements est cartésien, de sommet $X$ envoyé sur
$(x, f(x), gf(x))$, et la formule de projection donne
$\nu_{*}(1)$ pour ce plongement $\nu$ ; enfin $p_{13}\, \nu = \Gamma_{gf}$.

On obtient ainsi un foncteur $\varphi \colon \mathcal{V} \to \mathcal{V}^{A}$
qui est l'identité sur les objets.

Inversement, une correspondance agit : pour $u \in A(X \times Y)$ et
$x \in A(X)$, on pose
\[ \tilde{u}(x) = p_{2*}\bigl( u\, p_{1}^{*}(x) \bigr) \in A(Y) . \]
On a $\tilde{\Delta}_{X} = \mathrm{id}_{A(X)}$, par la formule de projection
appliquée à $\delta_{X}$ ; et
$\widetilde{v \circ u} = \tilde{v} \circ \tilde{u}$, les deux membres se
ramenant à $p_{Z*}\bigl(p_{YZ}^{*}(v)\, p_{XY}^{*}(u)\, p_{X}^{*}(x)\bigr)$ par
un échange dans un carré de projections et deux formules de projection. D'où
un foncteur $\psi \colon \mathcal{V}^{A} \to k\text{-Mod}$, $X \mapsto A(X)$,
$u \mapsto \tilde{u}$.

\pagerange{14}{16}
\subsection*{Le carré de foncteurs et le côté contravariant}

Le diagramme

\begin{tikzcd}
\mathcal{V} \arrow[r, "\varphi"] \arrow[d, "A"'] & \mathcal{V}^{A} \arrow[d, "\psi"] \\
\mathcal{C} \arrow[r, "\omega"'] & k\text{-Mod}
\end{tikzcd}

commute : pour $f \colon X \to Y$,
\[ \tilde{\gamma}_{f}(x) = p_{2*}\bigl(\Gamma_{f*}(1)\, p_{1}^{*}(x)\bigr) = p_{2*}\,\Gamma_{f*}\bigl(\Gamma_{f}^{*}\, p_{1}^{*}(x)\bigr) = f_{*}(x) , \]
puisque $p_{1}\Gamma_{f} = \mathrm{id}_{X}$ et
$p_{2}\Gamma_{f} = f$.\footnote{La page 14 termine le calcul par « $= x$ » ;
c'est $f_{*}(x)$ qu'il donne, et c'est ce qu'annonçait la ligne qui le
précède.}

Qu'en est-il de $f^{*}$ ? La symétrie $\sigma \colon Y \times X \to X \times Y$
définit la \emph{transposition}
${}^{t}u = \sigma^{*}(u) \in A(Y \times X)$ de $u \in A(X \times Y)$. Elle est
l'identité sur les objets et renverse la composition :
${}^{t}(v \circ u) = {}^{t}u \circ {}^{t}v$ ; c'est donc un isomorphisme de
catégories
\[ s \colon (\mathcal{V}^{A})^{\circ} \longrightarrow \mathcal{V}^{A} , \qquad s^{2} = \mathrm{id} . \]
Et l'image inverse est l'action de la transposée du graphe : pour
$f \colon X \to Y$ et $y \in A(Y)$,
\[ \widetilde{{}^{t}\gamma_{f}}\,(y) = f^{*}(y) , \]
par la même formule de projection. Autrement dit
$\varphi^{*} = s \circ \varphi^{\circ} \colon \mathcal{V}^{\circ} \to \mathcal{V}^{A}$
vérifie $\psi\, \varphi^{*} = \omega^{*} A^{\circ}$, où $\omega^{*}$ retient
$f^{*}$.\footnote{La page 14 écrit la question avec « $f \colon Y \to X$ », là
où $A(f)^{*} \colon A(Y) \to A(X)$ demande $f \colon X \to Y$, et note $s$
comme un foncteur de $\mathcal{V}^{A\circ}$ dans lui-même, alors que l'usage
qu'en fait la page 16 exige le but $\mathcal{V}^{A}$. La multiplicativité de
$s$ est marquée « vérifier » et l'identité $\psi\varphi^{*} = \omega^{*}A^{\circ}$
« à vérifier » ; toutes deux sont vraies sous l'hypothèse de signe. Sans elle,
${}^{t}(v \circ u) = (-1)^{|u||v|}\, {}^{t}u \circ {}^{t}v$. Une catégorie munie
d'une telle involution est ce qu'on appelle aujourd'hui une catégorie à
involution, ou catégorie « dague ».}

\subsection*{Formes et transposition}

Supposons $A(e) = k$.\footnote{« Sinon, on remplace $k$ par $A(e)$ », ajoute
la page 16.} L'unique flèche $\pi_{X} \colon X \to e$ donne une forme linéaire
$\kappa_{X} = \pi_{X*} \colon A(X) \to k$, et chaque $A(X)$ devient une algèbre
\emph{jaugée}.\footnote{La lecture « jaugés » est douteuse aux pages 16 et 18.
Le mot est certain page 41, où « algèbre jaugée » est définie comme une algèbre
munie d'une forme vers $k$ ; c'est le sens qu'on lui donne ici.} Pour
$u \in A(X \times Y)$, les applications $\tilde{u}$ et
$\widetilde{{}^{t}u}$ sont transposées l'une de l'autre pour les formes
bilinéaires $(a, b) \mapsto \kappa(ab)$ :
\[ \kappa_{Y}\bigl(y\, \tilde{u}(x)\bigr) = \kappa_{X \times Y}\bigl(p_{2}^{*}(y)\, u\, p_{1}^{*}(x)\bigr) = \kappa_{X}\bigl(\widetilde{{}^{t}u}(y)\, x\bigr) , \]
les deux égalités venant de la formule de projection et de
$\pi_{Y} p_{2} = \pi_{X} p_{1} = \pi_{X \times Y}$.\footnote{La page 16 écrit
l'hypothèse « $u$ central » et conclut « mod. signe ». Dans une algèbre
commutative au sens gradué, le signe est $(-1)^{|u||y|}$, trivial pour $u$ de
degré pair.} En particulier, si les $A(X)$ sont des algèbres de Poincaré, $A$
est connu dès qu'on connaît ces algèbres jaugées et les $f^{*}$ : la
transposition détermine les $f_{*}$. C'est ce qu'établiront les pages 33 et 35.

\pagerange{18}{18}
\subsection*{Ce que les calculs ont utilisé}

Aucune des vérifications précédentes n'a fait appel à l'addition des $A(X)$ :
seulement à leur multiplication. La théorie des correspondances pourrait donc
s'écrire pour des foncteurs à valeurs dans des « monoïdes bivariants ». Il
écarte aussitôt cette généralité pour une raison de signes : dès que les
formules de projection à gauche et à droite ne valent pas toutes deux sans
signe, il faudrait des monoïdes \emph{gradués} $(A^{i})_{i \in \mathbf{Z}}$,
avec $A^{i} \times A^{j} \to A^{i+j}$, munis d'involutions
$\sigma^{i} \colon A^{i} \to A^{i}$ triviales en degré pair, et d'un degré
$d(A)$ qui imposerait à $f_{*}$ le degré $d(B) - d(A)$. C'est la première
apparition du degré $d(A)$, qui gouvernera les pages 35 et 41 ; le projet des
monoïdes gradués, lui, n'est pas repris.\footnote{Le passage est biffé en bloc
sur la page 18 et plusieurs mots en sont illisibles, dont ce qu'il dit savoir
des signes à la dernière ligne ; « jaugés » y est une lecture douteuse, et
« gradués » conviendrait aussi bien au sens.}

\pagerange{20}{20}
\section*{Sous-foncteurs stables (pages 20 à 24)}

\subsection*{Le critère}

Pour tout $X$, soit $A'(X) \subset A(X)$ un sous-anneau (contenant $1$).
Les conditions suivantes sont équivalentes.

\begin{enumerate}
  \item Les $A'(X)$ sont stables par images inverses et par images directes
  par tous les morphismes de $\mathcal{V}$ (ils le sont par produit, étant des
  sous-anneaux).
  \item a) Pour $u \in A'(X \times Y)$ et $x \in A'(X)$, on a
  $\tilde{u}(x) \in A'(Y)$ ; b) pour $x \in A'(X)$, $y \in A'(Y)$, on a
  $x \boxtimes y = p_{1}^{*}(x)\, p_{2}^{*}(y) \in A'(X \times Y)$ ; c) pour
  tout $f$, $\gamma_{f} \in A'(X \times Y)$.
\end{enumerate}

\emph{Démonstration.} Si 1 vaut, a) résulte de la formule même de
$\tilde{u}$, b) de celle de $x \boxtimes y$, et c) de ce que $\gamma_{f}$ est
l'image directe de $1$. Si 2 vaut, $f_{*} = \tilde{\gamma}_{f}$ et
$f^{*} = \widetilde{{}^{t}\gamma_{f}}$ ; or ${}^{t}\gamma_{f} = \sigma^{*}\gamma_{f}$
est l'action du graphe de l'isomorphisme $\sigma^{-1}$, donc est dans $A'$ par
a) et c). La stabilité par produit intérieur suit de b) et de
$x y = \delta^{*}(x \boxtimes y)$.\footnote{La page 20 énonce la condition 1
aussi sous la forme « par produits $\boxtimes$ » et vérifie que les deux formes
s'équivalent, par $x \otimes y = \delta^{*}(x \boxtimes y)$ et
$x \boxtimes y = p_{1}^{*}(x)\, p_{2}^{*}(y)$. La stabilité par la symétrie,
qu'il invoque comme « stabilité par isom. », est déduite ici de a) et c).
Une note en marge demande un énoncé plus général, où « tout morphisme » serait
remplacé par une famille donnée de morphismes ; il n'est pas écrit.}

Il ajoute, sans démonstration, que si $\mathcal{V}$ a un objet final, ces
sous-foncteurs correspondent exactement aux sous-catégories de
$\mathcal{V}^{A}$ ayant les mêmes objets, stables par produit extérieur des
flèches, et contenant les morphismes de $\mathcal{V}$ et leurs transposés
--- les transposés des diagonales suffisant.\footnote{Une première version,
biffée, parlait de sous-catégories \emph{pleines}, ce qui est faux : une
sous-catégorie pleine ayant tous les objets est la catégorie entière. Le
passage d'un $A'$ à une sous-catégorie est
$\operatorname{Hom}(X, Y) = A'(X \times Y)$, et le retour
$A'(X) = \operatorname{Hom}(e, X)$ ; que ces deux opérations soient inverses
repose sur l'identification
$\operatorname{Hom}(X, Y) = \operatorname{Hom}(e, X \times Y)$, que réalisent
précisément les transposés des diagonales.}

\pagerange{22}{22}
\subsection*{Le plus petit sous-foncteur stable}

Une intersection de sous-foncteurs stables est stable ; il existe donc un plus
petit, $A_{\mathrm{sp}} \subset A$.\footnote{La page 22 dit « le plus petit
ss-foncteur stable engendré », et la ligne s'arrête : ce qui l'engendre n'est
pas écrit. Comme $1 \in A'(X)$ est imposé, le plus petit sous-foncteur stable
tout court est ce qu'on retient ici ; l'indice est le sien.} Il contient les
classes $f_{*}(1_{Y})$ pour tous les $f \colon Y \to X$, donc le sous-groupe
qu'elles engendrent. Quand $\mathcal{V}$ est la catégorie des variétés
projectives lisses sur un corps algébriquement clos de caractéristique $0$, et
$A$ une cohomologie comme celles de la page 41, ce sous-groupe est déjà stable
et $A_{\mathrm{sp}}(X)$ \emph{est} l'image de l'application classe de cycle
$\mathrm{CH}^{*}(X) \to A(X)$ : toute sous-variété fermée $Z \subset X$ admet
une résolution $\tilde{Z} \to X$ dont l'image directe de $1$ est la classe de
$Z$.\footnote{La page 22 dit que $A_{\mathrm{sp}}(X)$ est, « dans beaucoup de
cas importants », identique à l'ensemble des $f_{*}(1_{Y})$, en citant les
variétés projectives lisses et, en caractéristique $0$, la résolution des
singularités (démontrée par Hironaka en 1964). L'\emph{ensemble} ne suffit
pas : il n'est pas stable par produit, puisque sur l'éclaté d'un point dans
une surface le diviseur exceptionnel $E$ vérifie $E \cdot E = -[\mathrm{pt}]$,
qui n'est l'image directe de $1$ par aucun morphisme. C'est le
\emph{sous-groupe} engendré qui convient, et la page s'intitule d'ailleurs
« Sous-anneaux bivariants ». En caractéristique $p$, les altérations de de Jong
(1996) donnent le même énoncé à coefficients rationnels. Le mot « classes
algébriques » n'est pas sur la page.}

\pagerange{22}{24}
\subsection*{Engendrement}

Soit $A'$ stable. Une application
$A(X_{1}) \times \cdots \times A(X_{n}) \to A(X)$ est \emph{du type $A'$} si
elle s'écrit
\[ (x_{1}, \ldots, x_{n}) \longmapsto \tilde{u}(x_{1} \boxtimes \cdots \boxtimes x_{n}) , \qquad u \in A'(X_{1} \times \cdots \times X_{n} \times X) . \]
Les images directes et inverses sont du type $A'$ (par le critère), et les
composés d'applications du type $A'$ le sont encore : si $\varphi_{1}$,
$\varphi_{2}$ et $\chi$ le sont, $\chi(\varphi_{1}, \varphi_{2})$ l'est, la
correspondance qui la définit s'obtenant par composition et produit extérieur
de celles de $\chi$, $\varphi_{1}$, $\varphi_{2}$ ; le produit de deux
applications du type $A'$ l'est aussi.\footnote{Les pages 22 et 24 sont d'une
écriture rapide où la prose se perd entre les formules ; la définition, la
stabilité par composition et par produit s'y lisent, leur justification non.
L'argument donné ici est de nous. La page parle aussi d'applications
« $A'$-admissibles », mot qu'elle remplace par « du type $A'$ ».}

\emph{Application.} Donnons-nous, pour chaque $X$, une partie
$E_{X} \subset A(X)$, éventuellement vide. Le plus petit sous-foncteur stable
$B \subset A$ contenant $A'$ et les $E_{X}$ est donné par : $B(X)$ est
l'ensemble des sommes finies d'éléments
$\varphi(x_{1}, \ldots, x_{n})$, où $\varphi$ est du type $A'$ et
$x_{i} \in E_{X_{i}}$ ($n = 0$ donnant $A'(X)$). En effet ces $B(X)$ sont
stables par images directes et inverses et par produit, et tout sous-foncteur
stable contenant $A'$ et les $E_{X}$ les contient.\footnote{La page 24 dit
« l'ensemble des éléments de la forme » ; pour obtenir un sous-anneau il faut
les sommes finies, et on les ajoute. Le cas $n = 0$ est la mention
« $A'(X) \cup \ldots$ » de la page, dont la suite est illisible.} La page 50
fera exactement ce calcul pour $A' = A_{\mathrm{sp}}$ et pour la classe d'un
point.

\pagerange{26}{26}
\section*{Algèbres de Poincaré (pages 26 à 30)}

\subsection*{Une catégorie d'« espaces » algébriques}

Soit $\mathcal{P}$ la catégorie dont les objets sont les \emph{algèbres de
Poincaré} sur un corps $k$ : $k$-algèbres graduées commutatives au sens gradué,
de type fini, nulles en degré $< 0$, munies d'une forme linéaire
$\kappa_{A} \colon A \to k$ nulle hors du degré $d(A)$, telle que
$(x, y) \mapsto \kappa_{A}(xy)$ soit non dégénérée ; les morphismes sont les
homomorphismes d'algèbres graduées, sans condition sur les
formes.\footnote{« Strictement non dégénérée » est ajouté en interligne à la
page 26 ; sur un corps et en type fini, cela veut dire que la forme identifie
chaque $A^{i}$ au dual de $A^{d-i}$. La page 26 dit « de degré $d^{\circ}(A)$ »
et note cette catégorie $\mathcal{V}^{\circ}$. Ce sont les \emph{algèbres à
dualité de Poincaré}, qu'on appelle aussi algèbres de Frobenius graduées
commutatives. Une section du haut de la page, qui calculait le signe de la
symétrie, est barrée de deux traits obliques ; elle est reprise page 30.} On
lit $\mathcal{P}^{\circ}$ comme une catégorie d'espaces : un morphisme
$f \colon A \to B$ de $\mathcal{P}^{\circ}$ est un homomorphisme
$f^{*} \colon B \to A$, et on lui associe son \emph{adjoint}
$f_{*} \colon A \to B$, défini par
\[ \kappa_{B}\bigl(b\, f_{*}(a)\bigr) = \kappa_{A}\bigl(f^{*}(b)\, a\bigr) . \]
On vérifie que $f \mapsto f_{*}$ est fonctoriel, que la formule de projection
$f_{*}\bigl(a\, f^{*}(b)\bigr) = f_{*}(a)\, b$ est vraie (en degré pair), et
que $\kappa_{B} f_{*} = \kappa_{A}$. D'où un foncteur de
$\mathcal{P}^{\circ}$ dans les algèbres bivariantes graduées : c'est la
cohomologie d'un espace sans l'espace.

\pagerange{28}{28}
\subsection*{Produits et signes}

Dans $\mathcal{P}$, $k$ est objet initial et le produit tensoriel gradué
$A \otimes B$ est le coproduit ; c'est donc le produit de
$\mathcal{P}^{\circ}$, et $(f \times g)^{*} = f^{*} \otimes g^{*}$. On munit
$A \otimes B$ de la forme $\kappa_{A \otimes B}(a \otimes b) = \kappa_{A}(a)\,\kappa_{B}(b)$.

On s'attend à $(f \times g)_{*} = f_{*} \otimes g_{*}$. \emph{C'est faux} dès
que les degrés sont impairs. Avec la convention de signe de Koszul pour la
multiplication de $A \otimes B$ et $f_{*} \otimes g_{*}$ pris sans signe,
l'adjonction donne
\[ (f \times g)_{*}(a \otimes b) = (-1)^{|f|\,(d(B) + |b|)}\; f_{*}(a) \otimes g_{*}(b) , \qquad |f| = d(A') - d(A) , \]
pour $f \colon A \to A'$ et $g \colon B \to B'$ dans $\mathcal{P}^{\circ}$ : le
correctif est l'opérateur de signe du codegré sur $B$, élevé à la puissance
$|f|$. En particulier \emph{il n'y a aucun signe dès que $f$ est de degré
pair}.\footnote{La page 28 essaie successivement quatre conventions ---
forme tordue ou non, produit tensoriel d'applications tordu ou non --- et note
le signe parasite de chacune, pour conclure que la convention sans aucun signe
est la plus simple. La formule encadrée qui en résulte a la même forme que
celle-ci, $(f_{*} \otimes g_{*})$ composé avec
$\mathrm{id}_{A} \otimes [(-1)^{B}]^{\deg f}$, mais son indice et la
disposition de son second facteur se lisent mal, et plusieurs exposants des
lignes précédentes sont surchargés ; la formule est donc recalculée ici sous
les conventions dites, et sa conclusion --- aucun signe si $\deg f$ est pair
--- coïncide avec celle de la page. La règle de Koszul ne porte pas encore ce
nom sur la page.}

\pagerange{30}{30}
\subsection*{Degré pair, échange, et une question}

À partir de là, il se borne aux morphismes de degré pair, et
$(f \times g)_{*} = f_{*} \otimes g_{*}$. La symétrie
$s(a \otimes b) = (-1)^{|a||b|}\, b \otimes a$ de $A \otimes B$ sur
$B \otimes A$ vérifie alors, comme il convient,
$(g \times f)_{*} \circ s = s \circ (f \times g)_{*}$.

Il vérifie ensuite les deux formules d'échange pour le changement de base par
une projection. Si $R$ est une algèbre de Poincaré (un « espace » $S$), la
projection de $X \times S$ sur $X$ correspond à $p^{*}(a) = a \otimes 1$ et à
$p_{*} = \mathrm{id} \otimes \kappa_{R}$ ; pour $f' = f \times \mathrm{id}_{S}$
on a bien
\[ p^{*} f_{*} = f'_{*}\, q^{*} , \qquad f^{*} p_{*} = q_{*}\, f'^{*} , \]
la seconde parce que les deux membres envoient $a \otimes r$ sur
$\kappa_{R}(r)\, f^{*}(a)$. Il vérifie de même la formule d'échange pour le
carré de graphes de la page 10, calculée sur les éléments de
$A \otimes B \otimes C$.

Et il pose la question qui reste ouverte dans tout le dossier : trouver un cas
général où la formule d'échange est vraie --- suffit-il que le produit fibré
existe dans la catégorie et que les dimensions soient correctes ?\footnote{En
géométrie, la réponse est oui sous la forme de la transversalité : la formule
vaut quand la codimension du produit fibré est la somme des codimensions, et
elle est corrigée par une classe d'Euler sinon (formule d'intersection en
excès). Les pages 54 et 56 font ce décompte de dimensions pour les diagonales,
et y trouvent la caractéristique d'Euler ; le dossier ne relie pas ce calcul à
la question de la page 30.}

\pagerange{33}{33}
\section*{La catégorie des algèbres bivariantes (pages 33 et 35)}

Soit $k$ un anneau. La catégorie des \emph{algèbres bivariantes} a pour objets
les $k$-algèbres $A$ munies d'une forme $\kappa_{A} \colon A \to k$, et pour
morphismes $A \to B$ les couples $f = (f_{*}, f^{*})$ où
$f^{*} \colon B \to A$ est un homomorphisme d'algèbres unitaires et
$f_{*} \colon A \to B$ une application additive satisfaisant la formule de
projection et $\kappa_{B} f_{*} = \kappa_{A}$.\footnote{La page 33 écrit les
deux formes de la formule de projection, $f_{*}(x f^{*}(y)) = f_{*}(x) y$ et
$f_{*}(f^{*}(y) x) = y f_{*}(x)$, avec en marge « il faut choisir entre l'une et
l'autre formule ! » ; elles sont équivalentes dans le cas commutatif, et dans
le cas gradué commutatif si $f_{*}$ est de degré pair. Une note entre crochets
donne le signe $(-1)^{\delta j}$ qui les relie en général.}

Un tel morphisme est un isomorphisme si et seulement si $f^{*}$ est bijectif et
$f_{*}(1)$ inversible : la formule de projection donne
$f_{*}(f^{*}(y)) = y\, f_{*}(1)$, donc
$f_{*} = f_{*}(1) \cdot (f^{*})^{-1}$. Les « bons » isomorphismes sont ceux
où $f_{*}(1) = 1$, c'est-à-dire $f_{*} = (f^{*})^{-1}$ ; les autres existent,
par exemple $f^{*} = \mathrm{id}$, $f_{*}(y) = y\sigma$ avec $\sigma$ inversible
et $\kappa_{B}(y) = \kappa_{A}(y\sigma^{-1})$.\footnote{La remarque est écrite
en oblique dans la marge gauche de la page 33 ; l'exemple est celui de la
note (*) de la page 41, qu'on rapproche ici.}

\emph{L'image directe est déterminée par l'image inverse.} Soit
$\theta_{A} \colon A \to \check{A} = \operatorname{Hom}_{k}(A, k)$,
$\langle y, \theta_{A}(x) \rangle = \kappa_{A}(yx)$. Si $\theta_{B}$ est
injectif, tout morphisme $(f_{*}, f^{*})$ vérifie
\[ \kappa_{B}\bigl(y\, f_{*}(x)\bigr) = \kappa_{A}\bigl(f^{*}(y)\, x\bigr) , \]
ce qui détermine $f_{*}$ par $f^{*}$. Réciproquement, si
$f^{*} \colon B \to A$ est un homomorphisme unitaire dont la transposée envoie
$\operatorname{Im}\theta_{A}$ dans $\operatorname{Im}\theta_{B}$ --- ce qui
est automatique si $\theta_{B}$ est bijectif --- l'application $f_{*}$ définie
par cette formule satisfait la formule de projection (dans le cas commutatif,
ou gradué commutatif en degré pair) et $\kappa_{B} f_{*} = \kappa_{A}$ (faire
$y = 1$) ; $(f_{*}, f^{*})$ est alors un morphisme.

\pagerange{35}{35}
\emph{Le décalage des degrés.} Dans le cas gradué, supposons que $\kappa_{A}$
se factorise par la composante $A^{d(A)}$ et $\kappa_{B}$ par $B^{d(B)}$, et
$\theta_{B}$ injectif. Si $x$ est homogène de degré $i$, la forme
$y \mapsto \kappa_{A}(f^{*}(y)\, x)$ ne voit que les $y$ de degré $d(A) - i$ ;
la composante de $f_{*}(x)$ de degré $j$ a donc une image par $\theta_{B}$
nulle sauf si $d(B) - j = d(A) - i$, donc est nulle. Ainsi
\emph{$f_{*}$ est homogène de degré $d(B) - d(A)$}.

La catégorie de travail est donc celle des algèbres graduées commutatives au
sens gradué, munies d'un degré $d(A)$ et d'une forme nulle hors de
$A^{d(A)}$, les $f_{*}$ étant homogènes de degré $d(B) - d(A)$ ; avec, comme
sous-catégories naturelles, celle des algèbres commutatives de degré $d(A)$
pair, et celle des algèbres nulles hors de $[0, d(A)]$. Il y insiste :
\emph{il n'est pas question de se borner aux algèbres de Poincaré}, pas même en
cohomologie modulo torsion, parce que les constructions universelles en font
sortir.\footnote{La fin de la page 35 est en grande partie illisible ; on n'en
retient que la thèse, qui se lit, et le motif, qui se devine à « les
constructions universelles ». Les exemples de la page 43 (anneaux de Chow et
leurs quotients) montrent que les foncteurs qui l'intéressent ne sont pas à
valeurs dans les algèbres de Poincaré.}

\pagerange{37}{37}
\section*{Produit tensoriel et algèbres simpliciales (pages 37 et 39)}

Le produit tensoriel gradué est un foncteur sur la catégorie $\mathcal{C}$ des
algèbres bivariantes graduées, mais --- « malheureusement », dit-il --- ce
n'est pas un produit dans $\mathcal{C}$ : il manquerait un morphisme diagonal
$A \to A \otimes A$, c'est-à-dire une image directe
$\delta_{*} \colon A \to A \otimes A$ adjointe de la multiplication, et elle
n'existe pas en général. Elle existe quand $A$ et $B$ sont des algèbres de
Poincaré, et $A \otimes B$ est alors le produit dans la sous-catégorie
pleine qu'elles forment, puisque l'image directe y est déterminée par l'image
inverse et que $A \otimes B$ est le coproduit des algèbres.

Une \emph{algèbre bivariante simpliciale} est un foncteur contravariant
$I \mapsto A(I)$ de la catégorie des ensembles finis non vides, avec toutes
les applications, dans $\mathcal{C}$.\footnote{« Pas ½ simpliciale ! », écrit
la page 37 : toutes les applications, surjections comprises, et pas seulement
les injections. Indexé par les ensembles finis plutôt que par les ensembles
ordonnés $[n]$, c'est ce qu'on appelle aujourd'hui un objet simplicial
\emph{symétrique} ; le modèle est $I \mapsto H^{*}(X^{I})$.} Elle est
\emph{spéciale} si elle transforme sommes en produits tensoriels : pour
$I = I' \sqcup I''$, les images inverses par les deux inclusions définissent un
isomorphisme
\[ A(I') \otimes A(I'') \xrightarrow{\ \sim\ } A(I' \sqcup I'') , \]
ce qui est la formule de Künneth du modèle. Il affirme que, si l'on se borne
aux algèbres simpliciales spéciales dont la valeur sur un point est une
algèbre de Poincaré, $A \mapsto A(\mathrm{pt})$ est une \emph{équivalence} avec
la catégorie des algèbres de Poincaré.\footnote{Aucune démonstration n'est sur
la page, et la page 41, qui répète l'énoncé, le marque « à vérifier ». La
spécialité force $A(I) = A(\mathrm{pt})^{\otimes I}$, les images inverses par
les surjections sont alors les multiplications le long des fibres, et les
images directes sont déterminées par les images inverses (page 33) ; c'est ce
qui rend l'énoncé plausible. On remarque que la définition ne demande aucune
formule d'échange, et à raison : elle serait fausse pour le carré de deux
diagonales. La page écrit l'argument $A(\Delta(0))$, qui se lit mal.}

\pagerange{39}{39}
\subsection*{Une Proposition sur les systèmes stables}

Soit $X$ un espace dont la cohomologie satisfait Künneth et Poincaré. Pour
tout $n \geq 1$, soit $A_{n} \subset H^{*}(X^{n})$ une sous-algèbre stable par
le groupe symétrique. Les conditions suivantes sont équivalentes.

\begin{enumerate}
  \item[(i)] a) Pour $u \in A_{n+m} \subset H^{*}(X^{n} \times X^{m})$ et
  $x \in A_{n}$, on a $\tilde{u}(x) \in A_{m}$ ; b) $1 \in A_{n}$, et
  $x \in A_{n}$, $y \in A_{m}$ entraînent $x \boxtimes y \in A_{n+m}$ ; c) pour
  toute surjection $[1, n+1] \to [1, n]$, la classe du graphe de l'application
  correspondante $X^{n} \to X^{n+1}$ est dans $A_{2n+1}$, et de même pour les
  injections $[1, n] \to [1, n+1]$.
  \item[(ii)] Les $A_{n}$ sont stables par images inverses et directes par les
  projections $X^{n+1} \to X^{n}$ et par les plongements diagonaux
  $X^{n} \to X^{n+1}$.
  \item[(iii)] Vues comme ensembles d'homomorphismes $H^{*}(X^{n}) \to H^{*}(X^{m})$,
  les parties $A_{n+m}$ sont stables par composition et par $\boxtimes$,
  contiennent les identités et les images inverses et directes par les
  applications simpliciales.
\end{enumerate}

Ces conditions entraînent la stabilité par images inverses et directes pour
toutes les applications simpliciales $X^{n} \to X^{m}$.\footnote{La page 39
écrit en (i) a) « $u(x) \in H^{*}(X^{m})$ », ce qui ne dit rien ; c'est
$A_{m}$ qu'il faut. Elle ne donne pas de démonstration. C'est le critère de la
page 20 appliqué à la catégorie des puissances $X^{I}$ et des applications
$X^{\alpha}$ : ses produits sont $X^{I} \times X^{J} = X^{I \sqcup J}$, et
ses morphismes sont engendrés par les permutations, les projections et les
plongements diagonaux, d'où (i) $\Leftrightarrow$ (ii) ; (iii) est le NB de la
page 20 sur les sous-catégories. Le symbole de (iii) est un carré plein sur la
page, lu $\boxtimes$.}

La page 56 vérifiera ces conditions pour la sous-algèbre tautologique.

\pagerange{41}{41}
\section*{Seconde mouture : algèbres jaugées et correspondances entre puissances (pages 41 à 47)}

\subsection*{Les définitions reprises}

La page 41 reprend la catégorie des $k$-algèbres graduées bivariantes :
commutatives au sens gradué, munies d'un degré $d(A)$ et nulles hors de
$[0, d(A)]$ ; l'objet $k$, placé en degré $0$, en fait partie.\footnote{La page
lit l'intervalle $[1, d(A)]$, ce qui exclurait $k$ en degré $0$ mentionné à la
ligne suivante et l'unité de toute algèbre ; on corrige en $[0, d(A)]$.} Une
telle algèbre est \emph{jaugée} si elle est munie d'une forme vers $k$ ; dans
certains cas l'image directe est déterminée par l'image inverse (pages 33 et
35), mais même alors un isomorphisme $(f_{*}, f^{*})$ n'a pas besoin de
vérifier $f_{*} = (f^{*})^{-1}$, comme le montre l'exemple rapporté plus haut.

Pour les algèbres simpliciales, la spécialité demande maintenant deux choses :
l'isomorphisme de Künneth $A(I') \otimes A(I'') \simeq A(I' \sqcup I'')$ pour
les images inverses, et, pour $\alpha \colon I \to J$ et
$\beta \colon I' \to J'$, l'identification de
$(\alpha \sqcup \beta)_{*}$ à $\alpha_{*} \otimes \beta_{*}$ pour les images
directes. L'équivalence avec les algèbres de Poincaré est répétée. Il ajoute
que $A(\cdot) \mapsto \bigl(A(\mathrm{pt}), \Delta\bigr)$, où
$\Delta = \delta_{*}(1) \in [A \otimes A]^{d(A)}$ est la classe de la
diagonale, serait un foncteur pleinement fidèle vers les algèbres jaugées
munies d'un tel $\Delta$, en réservant les axiomes à imposer à
$\Delta$.\footnote{« À vérifier », en marge, avec la remarque qu'il faut
d'abord connaître $\delta_{*} \colon A \to A \otimes A$. La pleine fidélité
n'est claire que pour les objets spéciaux, où $A(I)$ est engendré par les
images de $A(\mathrm{pt})$ ; la page ne précise pas. Les axiomes réservés
sont, en termes d'aujourd'hui, ceux qui font de $\Delta$ l'inverse de la forme
$\kappa(xy)$ --- les identités « en zigzag » d'une algèbre de Frobenius, où
$\delta_{*}$ est la comultiplication. Qu'une algèbre de Frobenius commutative
soit la même chose qu'une théorie topologique des champs en dimension $2$ est
un théorème d'Abrams (1996) ; rien de tel n'est sur la page.}

\pagerange{41}{43}
\subsection*{Trois exemples}

\begin{enumerate}
  \item Une cohomologie sur la catégorie $\mathrm{Var}$ des variétés complètes
  (lisses), satisfaisant la formule de Künneth et la dualité de Poincaré avec
  un isomorphisme de trace en dimension maximale, est un foncteur vers les
  algèbres graduées bivariantes jaugées ; appliqué à l'objet simplicial
  $I \mapsto X^{I}$, il donne une algèbre bivariante jaugée simpliciale
  spéciale, déterminée par $H^{*}(X)$. Variantes : variétés topologiques
  compactes ; coefficients dans un anneau, à condition que la dualité de
  Poincaré vaille pas seulement modulo torsion.\footnote{C'est à peu près ce
  qu'on appelle depuis Kleiman (1968) une cohomologie de Weil, moins
  l'application classe de cycle. Le mot que la transcription lit
  « trivialisation », pour l'isomorphisme de trace, est une lecture d'après
  l'argument.}
  \item L'anneau de Chow $A^{*}(X)$, jaugé par le degré des $0$-cycles, ou ses
  quotients par une relation d'équivalence : numérique, homologique,
  « $\tau$-algébrique », et les quotients de Picard et d'Albanese. On obtient
  une algèbre bivariante jaugée simpliciale qui « n'a plus aucune tendance à
  être spéciale », et dont les composantes ne sont pas de Poincaré ; on peut
  tensoriser par $\mathbf{Q}$.\footnote{Le $\tau$ est ajouté en interligne ;
  la page ne définit pas cette équivalence. Le mot est proche de la
  $\tau$-équivalence de SGA 6 (équivalence algébrique à torsion près), sans que
  rien ne permette de l'identifier. Ce sont les « relations d'équivalence
  adéquates » de la théorie des motifs purs.}
  \item La $K$-théorie $K(X)$, jaugée par la caractéristique d'Euler-Poincaré
  $\chi$ et considérée comme graduée en degré $0$.
\end{enumerate}

\pagerange{43}{45}
\subsection*{Correspondances entre puissances}

Soit $c \in A(I \sqcup I')$, et $\alpha \colon I \to I \sqcup I'$,
$\alpha' \colon I' \to I \sqcup I'$ les inclusions (qui correspondent aux
projections de $X^{I \sqcup I'}$ sur $X^{I}$ et $X^{I'}$). La correspondance
$c$ agit par
\[ \tilde{c} \colon A(I) \xrightarrow{\ \alpha^{*}\ } A(I \sqcup I') \xrightarrow{\ y \mapsto c\,y\ } A(I \sqcup I') \xrightarrow{\ \alpha'_{*}\ } A(I') , \]
avec, note-t-il, un choix entre $cy$ et $yc$. Ces opérations ne demandent que
les lois d'anneau et les $\alpha^{*}$, $\alpha_{*}$ pour les injections, et même
seulement pour les « opérations face », où $J$ a un élément de plus que $I$.

Écrivant $A(I', I)$ pour $A(I \sqcup I')$, correspondances de $I$ vers $I'$, il
définit la composition
\[ A(I'', I') \times A(I', I) \to A(I'', I) , \qquad v \circ u = p_{13*}\bigl(p_{12}^{*}(u)\, p_{23}^{*}(v)\bigr) , \]
les projections étant celles de $X^{I \sqcup I' \sqcup I''}$, et vérifie
$\widetilde{v \circ u} = \tilde{v} \circ \tilde{u}$ : les deux membres valent
$p_{3*}\bigl(p_{12}^{*}(u)\, p_{23}^{*}(v)\, p_{1}^{*}(x)\bigr)$.\footnote{L'ordre
des facteurs est l'inverse de celui de la page 2 ; les deux compositions
diffèrent du signe $(-1)^{|u||v|}$ et coïncident sous l'hypothèse de signe. La
page trace la lettre de $p_{32}^{*}(\cdot)$ comme un $w$ ; la formule demande
$v$. Au dernier pas de la vérification, les deux produits apparaissent dans
des ordres différents, d'où la marge « OK mod signes ; faire attention ».}
Le calcul de $\tilde{v}(\tilde{u}(x))$ passe par l'échange dans le carré
\[ X^{I \sqcup I' \sqcup I''} \to X^{I' \sqcup I''} , \qquad X^{I \sqcup I'} \to X^{I'} , \]
et c'est là qu'il note en marge que l'« axiome cartésien » a été oublié dans
les définitions de cette mouture.

\pagerange{47}{47}
\subsection*{Unités, et l'arrêt}

L'application codiagonale $I \sqcup I \to I$ correspond à la diagonale
$\delta \colon X^{I} \to X^{I \sqcup I}$ ; on pose
$\Delta_{I} = \delta_{*}(1) \in A(I \sqcup I)$, et
$\tilde{\Delta}_{I} = \mathrm{id}_{A(I)}$ par la formule de projection, comme
à la page 12.\footnote{La page écrit $\tilde{\delta}_{I}$ au moment de
l'énoncé, $\Delta_{I}$ dans la définition.} Il annonce qu'il reste à vérifier
l'associativité de la composition et que les $\Delta_{I}$ sont des unités ---
ce que les pages 4 et 6 avaient fait dans la première mouture --- et la page
s'arrête au milieu d'une phrase.\footnote{« Il faut encore \emph{définir}
l'associativité » : le verbe est douteux, et le sens demande
« vérifier ». La page 47 se termine sur une virgule ; la suite n'est pas dans
le dossier.}

\pagerange{50}{50}
\section*{La sous-algèbre tautologique de $H^{*}(X^{n})$ (pages 50 à 56)}

\emph{Cadre.} Soit $X$ une variété projective lisse connexe sur $\mathbf{C}$,
de dimension réelle $n$ (donc $n$ pair), et $H^{*}$ sa cohomologie rationnelle
--- ou toute cohomologie de Weil ; soit $y \in H^{n}(X)$ la classe d'un point
et $\chi$ la caractéristique d'Euler de $X$.\footnote{La page 50 ne dit rien
de $X$ ; elle écrit seulement que $y$ est de degré $n$, que $y^{2} = 0$, et que
$\chi \in \mathbf{Z}$ est « la car. d'EP ». Ces hypothèses sont celles qui
rendent vraies ses relations. Pour $n$ impair, $\chi = 0$ et les signes
compliquent les formules.} Pour un ensemble fini $I$, on note $y_{i}$
l'image inverse de $y$ par la projection d'indice $i$, et $\delta_{I}$ la
classe de la petite diagonale de $X^{I}$, de degré $(|I| - 1)\,n$. La
\emph{sous-algèbre tautologique} $T(X^{I}) \subset H^{*}(X^{I})$ est la
sous-algèbre engendrée par les $y_{i}$ et les classes des diagonales partielles.

\subsection*{Les premières tables}

$T(X)$ a pour base $1$, $y$, avec $y^{2} = 0$.

$T(X^{2})$ est engendrée par $y_{1}$, $y_{2}$ et $\delta = \delta_{\{1,2\}}$,
de degré $n$, avec les relations
\[ y_{1}^{2} = y_{2}^{2} = 0 , \qquad y_{1}\delta = y_{2}\delta = y_{1}y_{2} , \qquad \delta^{2} = \chi\, y_{1}y_{2} ; \]
elle est engendrée comme espace vectoriel par $1$ ; $y_{1}, y_{2}, \delta$ ;
$y_{1}y_{2}$.\footnote{La page écrit « $\delta \in H^{n}(X)$ » ; la diagonale
vit dans $H^{n}(X \times X)$. Elle appelle « base » ces familles ; ce sont des
familles génératrices, libres seulement en général : pour la sphère $S^{2}$,
où $H^{*}(X)$ se réduit à $k \oplus k y$, on a $\delta = y_{1} + y_{2}$. La
relation $\delta^{2} = \chi\, y_{1}y_{2}$ est la formule de Lefschetz pour
l'identité : l'auto-intersection de la diagonale compte les points fixes.}

$T(X^{3})$ est engendrée en degré $n$ par $y_{1}, y_{2}, y_{3}$ et
$\delta_{12}, \delta_{23}, \delta_{13}$, en degré $2n$ par la petite
diagonale $\delta = \delta_{\{1,2,3\}}$, avec les relations
\[
\begin{gathered}
y_{i}^{2} = 0 , \qquad \delta_{ij}^{2} = \chi\, y_{i}y_{j} , \qquad \delta_{12}\delta_{23} = \delta_{12}\delta_{13} = \delta_{23}\delta_{13} = \delta , \\
\delta_{ij}\, y_{i} = \delta_{ij}\, y_{j} = y_{i}y_{j} , \qquad \delta\, y_{i} = y_{1}y_{2}y_{3} , \qquad \delta^{2} = 0 , \qquad \delta_{ij}\,\delta = \chi\, y_{1}y_{2}y_{3} ,
\end{gathered}
\]
et comme famille génératrice : $1$ ; les six classes de degré $n$ ;
$y_{1}y_{2}, y_{2}y_{3}, y_{1}y_{3}, \delta_{12}y_{3}, \delta_{23}y_{1}, \delta_{13}y_{2}, \delta$
en degré $2n$ ; $y_{1}y_{2}y_{3}$ en degré $3n$.

\pagerange{52}{52}
\subsection*{Générateurs de $T(X^{k})$}

En degré $n$ : les $k$ classes $y_{i}$ et les $\binom{k}{2}$ diagonales
$\delta_{ij}$ ; en degré $rn$, pour $1 \leq r \leq k-1$ : les
$\binom{k}{r+1}$ petites diagonales $\delta_{I}$ des parties $I$ à $r+1$
éléments, jusqu'à la petite diagonale de $X^{k}$ en degré
$(k-1)n$.\footnote{La page note $\delta_{r}$ la petite diagonale de $X^{r+1}$
et ses « transformés » par les permutations ; c'est $\delta_{I}$ ici. Comme
générateurs d'algèbre, les $\delta_{I}$ avec $|I| \geq 3$ sont redondants :
$\delta_{\{1, \ldots, r+1\}} = \delta_{12}\,\delta_{23} \cdots \delta_{r,r+1}$,
les diagonales partielles se coupant transversalement. La page s'arrête
ensuite sur le mot « relations ».}

\pagerange{54}{56}
\subsection*{Une famille génératrice et sa table de multiplication}

Pour décrire $T(X^{I})$ comme espace vectoriel, il indexe ses classes par des
paires $(R, J)$ : $R$ une partition de $I$ en parties non vides et $J$ une de
ses parties, ou la partie vide.\footnote{La page 54 obtient la partie vide en
convenant que $\emptyset \in R$ ; c'est la même chose.} Pour un point
$a \in X$, soit
\[ Z_{R,J} = \bigl\{ (x_{i})_{i \in I} \;:\; x_{i} = a \text{ si } i \in J , \ x_{i} = x_{j} \text{ si } i, j \text{ sont dans une même partie de } R \bigr\} , \]
et $\varepsilon_{R,J}$ sa classe. À l'ordre des facteurs près,
\[ \varepsilon_{R,J} = \Bigl(\,\boxtimes_{i \in J}\, y\Bigr) \boxtimes \Bigl(\,\boxtimes_{P \in R,\ P \neq J}\, \delta_{P}\Bigr) , \]
de sorte que les $\varepsilon_{R,J}$ engendrent $T(X^{I})$ comme espace
vectoriel dès qu'on sait qu'elles en forment une sous-algèbre --- ce que
donne leur table de multiplication.

Soient $(R, J)$ et $(R', J')$, et $R''$ la plus fine partition moins fine que
$R$ et $R'$. Pour une partie $B$ de $R''$, notons $b$ et $b'$ les nombres de
parties de $R$ et de $R'$ contenues dans $B$ ; comme $B$ est « connexe » pour
$R$ et $R'$, on a toujours $b + b' \leq |B| + 1$.

\textbf{Proposition.} Le produit $\varepsilon_{R,J}\,\varepsilon_{R',J'}$ est
nul si $J$ et $J'$ sont non vides et contenus dans une même partie de $R''$
(en particulier si $J \cap J' \neq \emptyset$). Sinon, il est le produit
extérieur, sur les parties $B$ de $R''$, des classes $c_{B} \in H^{*}(X^{B})$
suivantes :
\begin{itemize}
  \item si $B$ contient $J$ ou $J'$ : $c_{B} = \prod_{i \in B} y_{i}$ si
  $b + b' = |B| + 1$, et $c_{B} = 0$ sinon ;
  \item si $B$ ne contient ni $J$ ni $J'$ : $c_{B} = \delta_{B}$ si
  $b + b' = |B| + 1$ ; $c_{B} = \chi \prod_{i \in B} y_{i}$ si
  $b + b' = |B|$ ; $c_{B} = 0$ si $b + b' < |B|$.
\end{itemize}

\emph{Démonstration.} Ensemblistement, et pour deux points $a \neq a'$,
$Z_{R,J} \cap Z_{R',J'}$ est vide dans le premier cas, et sinon le produit,
sur les parties de $R''$, d'un point ou d'une petite diagonale. On se ramène
ainsi à une seule partie $B$, puis au cas où $J = J' = \emptyset$ et où $R''$
est la partition grossière : c'est la réduction de la page 54. Posons
$k = |B|$, $s = b$, $s' = b'$. Alors $Z_{R}$ est isomorphe à $X^{s}$, de
codimension $(k - s)n$ dans $X^{k}$, de même pour $Z_{R'}$, et leur
intersection est la petite diagonale, de codimension $(k-1)n$. Comme la
codimension d'une intersection ne dépasse pas la somme des codimensions,
$s + s' \leq k + 1$. S'il y a égalité, l'intersection est transverse et le
produit est $\delta_{B}$. Si $s + s' < k$, le produit est de degré
$(2k - s - s')n > kn$, donc nul. Si $s + s' = k$, le produit est en degré
maximal ; l'intersection est de dimension $n$ au lieu de $0$, le fibré d'excès
est le fibré tangent de $X$ restreint à la diagonale, et la formule
d'intersection en excès donne $\chi$ fois la classe d'un point. Pour une partie
contenant $J$, le même décompte, où l'intersection attendue est un point, ne
laisse que le cas transverse.\footnote{La page 54 donne la nullité si
$J \cap J' \neq \emptyset$ et si les parties de $R''$ engendrées par $J$ et
$J'$ coïncident, une formule générale dont les exposants et les indices sont
surchargés au point de ne plus se lire, et la réduction à la partition
grossière. La page 56 fait le décompte des codimensions et encadre
$s + s' \leq k + 1$ ; elle trouve la diagonale dans le cas d'égalité, la
nullité « pour raison de degrés » si $k > s + s'$, et laisse le cas
$k = s + s'$ avec un $\chi$ biffé suivi d'un point d'interrogation. Le
coefficient $\chi$ est rétabli ici ; il est confirmé par deux relations de la
page 50, $\delta^{2} = \chi\, y_{1}y_{2}$ ($k = 2$, $s = s' = 1$) et
$\delta_{12}\,\delta = \chi\, y_{1}y_{2}y_{3}$ ($k = 3$, $s = 2$, $s' = 1$).
Le traitement des parties marquées d'un point, et l'identification du fibré
d'excès, sont de nous. La page désigne par $\delta_{k}$ la petite diagonale de
$X^{k}$, que sa propre page 52 noterait $\delta_{k-1}$. Les losanges d'algèbre
linéaire dessinés au bas de la page 56, avec des codimensions $p$, $q$, $r$,
sont le décompte correspondant pour des sous-espaces vectoriels ; leurs
étiquettes se lisent en partie seulement.}

\subsection*{Stabilités, et la conclusion}

Le sous-espace engendré par les $\varepsilon_{R,J}$ est donc une sous-algèbre ;
il est stable par $\boxtimes$ (clair), par images inverses par les projections
(on ajoute une partie à un élément), et par images directes par la projection
$X^{I} \to X^{I \smallsetminus \{i\}}$ : si $i \in J$, $Z_{R,J}$ s'envoie
isomorphiquement sur $Z_{R, J \smallsetminus \{i\}}$ ; si $i \notin J$ et que
la partie $P$ contenant $i$ a au moins deux éléments, $\delta_{P}$ s'envoie sur
$\delta_{P \smallsetminus \{i\}}$ ; si $P = \{i\}$, l'image directe est
nulle.\footnote{La page 56 écrit « on trouve $\delta_{I}$ », avec un trait sur
l'indice ; pour des raisons de degré c'est $\delta_{I \smallsetminus \{i\}}$.
Un d) qui suivait, avec trois lignes dont une encadrée, est biffé et
surchargé.} Les images inverses et directes par les plongements diagonaux
$\delta \colon X^{I \smallsetminus \{j\}} \to X^{I}$ s'y ramènent :
$\delta_{*}(z) = \delta_{ij}\, p^{*}(z)$ et
$\delta^{*}(z) = p_{*}(\delta_{ij}\, z)$, où $p$ est la projection qui oublie
$j$. Le système des $T(X^{n})$ satisfait donc les conditions de la
Proposition de la page 39, et il est stable par composition des
correspondances.\footnote{Les relations $\delta_{*}(z) = \delta_{ij}\,p^{*}(z)$
et $\delta^{*}(z) = p_{*}(\delta_{ij}\,z)$ découlent de la formule de
projection et de $p\,\delta = \mathrm{id}$ ; elles ne sont pas sur la page, qui
écrit seulement quelques flèches entre diagonales, $\Delta_{p} \leftarrow
\Delta_{q}$ et $\Delta_{p} \sqcup 1 \leftarrow \Delta_{q} \sqcup 1$, puis
« Stable par comp. ».}

La conclusion est encadrée : \emph{la sous-algèbre tautologique est
engendrée, par produits extérieurs, par la classe $y$ et les petites
diagonales $\delta_{I}$}. C'est donc le plus petit système stable, au sens des
pages 24 et 39, qui contienne la classe d'un point : chaque $\delta_{I}$ est
l'image directe de $1$ par le plongement de la petite diagonale, et y est
forcé.\footnote{Ce rapprochement avec la construction de la page 24 est de
nous. La page 56 se termine sur « $(\delta_{2} \boxtimes \delta_{2})$ », isolé.
Le sous-anneau de $\mathrm{CH}^{*}(X^{n})$ engendré par les diagonales et
quelques classes de base a été étudié récemment, pour les surfaces K3, par
Beauville et Voisin (2004) puis par Voisin ; on l'y appelle aussi
tautologique. Le dossier travaille en cohomologie, où la question est plus
simple.}

\end{document}
